\chapter{落地之后}
\chapterkana{降下の後}
\epigraphbox{“去者日以疏，生者日以亲。”}{——《古诗十九首·去者日以疏》}

\kaishou{十二万人落进黄河流域，二百一十七人借着另一批船往东渡海。\par{}不过四代，原本记得月亮的那群人便只剩百分之四点二。\par{}再往后，「华夏神族」四个字又在正史与谱牒里空挂了两千年。}

\section{四代衰减与记忆消退}

对于从天上下来的人而言，公元前二二一年的落地是一场被迫终止的迁徙，谈不上什么凯旋。双脚一踏进中原地面，等于重新回到那个生死流转的世界，两万年来被净秽系统隔绝在外的一切便立刻扑了上来。依据第 117 全宗第 96 叶与《广寒宫残简》简二十六保存的一组简略记录，落地之后大约立刻同时恢复了三件事。

第一，生理变化恢复，不再保持不变的躯体重新开始衰老与死亡；第二，生育恢复，初代之后重新诞下会老亦会死的普通凡人；第三，原本在无秽状态下不起作用的随身记忆消除场重新开始发作，他们在逐年衰老的同时，也将月面上的经历与整套精密技艺逐代淡忘。\footnote{关于落地后随身之场的复发、文献空白与代际遗忘的具体机制，参见索科洛夫《记忆消除场的文献学后果》；本书只照档案数字记录其可见后果，不对场的物理机制加以推演。}

如果将十二万落地者记作第一代，每隔大约五十年算作一代，先后四代合计不过二百年，手头保存下来的世代存活率数字一路跌落得惊人（见图 \ref{fig:survival}）。以十二万初代为分母计算，第一代尚存百分之六十一点四，第二代跌至百分之三十八点七，第三代降到百分之十二点九，及至第四代，留在世上者已不足最初人数的百分之四点二，折算下来大约只有五千人。\footnote{四项百分比见《月之都——〈广寒宫残简〉辑佚与考释》简二十六与第 117 全宗第 96 叶，分母均为 120{,}000。简记四代以后「无记」，第 117 全宗第 97 叶末行仅存铅笔字一行：「他们开始挖洞。」}

\begin{figure}[htbp]
  \centering
  \begin{tikzpicture}[scale=1.0, font=\footnotesize]
    \draw[line width=0.5pt] (0,0) -- (9.6,0);
    \draw[line width=0.5pt] (0,0) -- (0,4.6);
    \foreach \y/\lbl in {0/0\%, 1.0/20\%, 2.0/40\%, 3.0/60\%, 4.0/80\%}{
      \draw[line width=0.3pt] (-0.12,\y) -- (0,\y);
      \draw[line width=0.2pt, black!15] (0,\y) -- (9.6,\y);
      \node[anchor=east] at (-0.16,\y) {\lbl};
    }
    \node[anchor=south, rotate=90] at (-1.05,2.2) {占初代（十二万）比例};
    \fill[black!55] (1.1,0) rectangle (2.1,3.07);
    \fill[black!55] (3.3,0) rectangle (4.3,1.935);
    \fill[black!55] (5.5,0) rectangle (6.5,0.645);
    \fill[black!55] (7.7,0) rectangle (8.7,0.21);
    \node[anchor=south] at (1.6,3.08) {61.4\%};
    \node[anchor=south] at (3.8,1.945) {38.7\%};
    \node[anchor=south] at (6.0,0.655) {12.9\%};
    \node[anchor=south] at (8.2,0.22) {4.2\%};
    \node[anchor=north] at (1.6,-0.15) {第一代};
    \node[anchor=north] at (3.8,-0.15) {第二代};
    \node[anchor=north] at (6.0,-0.15) {第三代};
    \node[anchor=north] at (8.2,-0.15) {第四代};
    \node[anchor=north, font=\scriptsize, color=black!65] at (4.8,-0.75) {世代间隔约五十年（共约二百年，至第四代约余五千人）};
  \end{tikzpicture}
  \caption{落地后四代人口存活率递减（以初代十二万人为基准）}
  \label{fig:survival}
  \figsource{据俄罗斯科学院第 117 全宗第 96 叶及《月之都——〈广寒宫残简〉辑佚与考释》简二十六绘制。}
\end{figure}

两百年不到，作为一个独立识别群体的月面来客便在认同层面上消散殆尽：第四代之后的后裔不再觉得月亮与自己有何干系，也不再自称来自天上。至于血统本身，则在中原漫长的通婚混居里又延续了一千余年，直到唐宋以后彻底与周边原住民混作一团，再也分不出彼此。

\section{前二一〇年东渡与「天文密葬法」}

不过在这场大衰减刚刚开了个头的公元前二一〇年，即最后一批落地十一年之后、东郡坠星次年，有一小支月之民判断中原难以久留，决定再次渡海离去。他们动身的理由十分现实，摆在眼前的有三条：第一，秦始皇此时正四处遣使求取长生不死之药，而刚刚从天上下来、尚未在短时间内尽数老去的他们自己正是现成的寻药靶子；第二，撤离前早已测算出来的四代衰减曲线正在眼皮底下逐年兑现；第三，沾染污秽之后，连这群原本在城里不老的人也切身感受到了衰老与记忆散失的逼近。

借谁的船出海，史书上正好留着一笔对得上的账。《史记·淮南衡山列传》记方士徐福屡次上书秦廷，声称海中有三神山，请得「童男女三千人」「五谷」「百工」出海求仙。\footnote{《史记》卷一一八《淮南衡山列传》，中华书局 1959 年点校本，第 3086 页。该传记徐福得平原广泽而止，王之不来，结构本类移民而非寻药方队。}一支连童男女带粮食和各行工匠一齐配齐的船队，并不像去海上采药，倒像一次正经的渡海移民。东渡残档与后世名籍核对，共有二百一十七名月之民混进了这支船队，领头的是八意永琳与蓬莱山辉夜两个人。需要专门交代一句：这一支东渡者谁也没有服过蓬莱药；无论是八意永琳还是蓬莱山辉夜，与其余两百多人一样都是普通的月之民，一落地面便已沾染污秽，踏上东渡船只的时候照样在一点一点变老。后世传说将她们改写成服过不死药后升天返月，不过是把「再也回不去那座空城」拧成了一个比较体面的结局。

八意永琳登上徐福船的时候，身边带走的东西极少，概括起来只有两样：后羿那桩旧案的全部档册，以及她本人脑子里记着的「天文密葬法」。第七章留到此处方才交待全貌的那个问题，到这里总算能够闭合：为什么档案要把夹在《历一》里的那段隐秘算法叫作「天文密葬法」。这门技法的真正创制者正是八意永琳自己，与后羿没有关系；所谓「密葬」，是指以精密星位推演把原本在天体轨道上运行的目标掩蔽、挪移甚至葬入不可见处之术，也就是第七章第四十一枚所记「以术正之」的整套算理。这套技法当年本来就只有极少几个人掌握，无法像器物那样装箱搬运，能跟着她离开中原的，只有她脑海里的记忆本身。可惜人既落了地、沾了地面的污秽，原本失效的记忆消除场便迟早会替她把这门技法慢慢磨去，带走也就等于带去一块终究要融化掉的冰。

\section{九百年与天津神}

公元前二一〇年东渡的这二百一十七人，大致在弥生时代中期登上了日本列岛。等到当地真正用文字把这段古史定格成书，已经是公元七一二年的《古事记》与七二〇年的《日本书纪》，两书只隔八年，距离前二一〇年落地却隔着整整九百二十多年，中间既没有同期汉文实录，也缺乏连续不断的铭文校核。月之民讲给自己子孙听的那些城中往事，便在这九百年里一点一点溶进了弥生与古坟时代各部族的口传故事里，最后沉淀成「天津神」那一支神系。

在这里必须先划清一条界线：日后记纪里地位最显赫的天照大神、素盏鸣尊诸神，属于日本列岛原有的本土神话体系，跟这二百一十七个东渡者没有血缘或渊源上的联系；真正能够跟月面档案严丝合缝对上的，只有「月读」及「思兼」那寥寥数支。同时，第九章曾暂且悬置「被中原称为西王母的那位月面传历者在档案中究竟叫作何名」这一疑问，如今合参东渡名录与日本口传材料，便能正式作答：第九章表格里的「西王母／传历者」，与前二一〇年率众东渡并携走《天文密葬法》的八意永琳，实为同一人。现将全书涉及的跨地域个人名号与群体称谓并排收入表 \ref{tab:names}。

\begin{table}[htbp]
  \centering
  \caption{名号与群体称谓跨地域对照}
  \label{tab:names}
  \footnotesize
  \begin{tabular}{@{}p{25mm}p{26mm}p{32mm}p{39mm}@{}}
    \toprule
    档案本位／身份 & 中原文献称谓 & 日本文献称谓 & 说明 \\
    \midrule
    月之民（全体） & 华夏神族（12 万） & 天津神（217 人支系） & 群体统称，后世各自另起名号 \\
    地之主任 & 姮娥／嫦娥 & — & 避汉文帝刘恒讳改「嫦娥」，留于月面 \\
    城主／造船首脑 & — & 月读命（即月夜见） & 约前 1980 年死于月面，随口传抵日本 \\
    上级／历法主修者 & 西王母 & 八意思兼神（八意永琳） & 曾视察夏地实验，前 210 年东渡 \\
    东渡同行者 & — & 赫映姬／栉笥名（辉夜） & 《竹取物语》素材原型，同船未返月 \\
    \bottomrule
  \end{tabular}
  \tabsource{对照《古事记》《日本书纪》《竹取物语》《山海经》与 \texttt{ФВ-1970-18}；天照、素盏鸣等本土神系不入此表。}
\end{table}

拿表 \ref{tab:names} 里的三条对应细看，就能摸出跨海口传和本土改编各自的规律。第一条是月夜见对应月读命：他本人早在公元前一九八〇年就在净秽山里死在了月面上，根本没有乘过徐福的船，但因为他在月之民记忆里的地位极高，名号随着二百一十七人的口头追述传到列岛，九百年后成了三贵子里的月神。第二条是八意永琳对应八意思兼神，「思兼」本来就是思虑兼备之义，《古事记》里开天岩户的计谋正由他设局献策，保留了她原本作为制度与算法设计者的影子；而她在公元前二〇〇〇年前后以「西王母」之名下地视察历法实验时，留给夏人的印象在《山海经》里还是豹尾虎齿的荒远之神，即便她本人已经在前二一〇年离开中原，西王母的名号依然在华夏土壤里被人一改再改，到六朝《汉武帝内传》时期改成了端庄华贵的绝世女仙。

第三条则是蓬莱山辉夜与《竹取物语》里的赫映姬。\footnote{《竹取物语》现存文本成书约在九世纪末至十世纪初，晚于东渡千年余；定位仍归文学性口传整理，不因晚出即一笔抹杀，亦不作实录直接照搬。}东渡之时她同永琳一样未服过药，同列岛上普通居民般慢慢变老、老去，而那座留在轨道上的月之都早就在前二二一年断水断气、人去城空。只是在后世民间故事的讲法里，群仙之乡总比一座空壳更有说头，「回不去」被改写成「拒绝凡尘、重返天界」，于是才多出了那句辞别养父母、披上羽衣升天的收尾。

\section{名号的攀附}

海那边的二百一十七人渐渐溶进了弥生人群的口传谱系，留在中原的十二万人留下的却是另一种截然不同的东西。他们的血统早就在四代以后淡得看不见了，「华夏神族」这四个字却在正史、谱牒和诏敕里整整挂了两千多年。

读到这里，最容易犯的错误便是把「血统」和「名号」当成一回事。前面图 \ref{fig:survival} 算得十分清楚，落地四代、不过两百年的光景，原本自认为来自月面的一群人就已经缩减到区区数千人，到了唐宋以后更是与本土土著再难分清彼此。可就在这个血统已经稀得几乎找不到的漫长时期里，地面上自称「华夏神族」或其后裔的人反而越冒越多。

他们争着认这四个字，原因既不神秘，也不浪漫，无非是因为这个名号好用。按照月之都降落前后带下来的底子，这批人的第一代和二三代手里捏着比本土更准一点的历数、度量和营建规矩；前二二一年秦始皇之所以能一举把全天下的度量衡统一起来（与第三章考过的小篆、一里四百一十五点八米、高台院落全套对得上），背后正少不了这批人手上那几把旧尺和历算资格的影子。往后几百年，哪一家想办改正朔、行封禅，或者在士族郡望里抬高门第，谱系里挂上这一支，说话就能压住旁人一头。这情形与当年欧洲日耳曼各部抢着自称罗马继承人颇为相似：真正管用的向来是那块大家公认好使的招牌；至于血管里还剩多少先祖的余温，反倒没多少人去核实。

如果按朝代把这套名号的用法切开来看，前后刚好分成三段（见表 \ref{tab:lineage}）。各段之间顺着官修谱牒和改朝换代的需要逐次变样，起落分明。

\begin{table}[htbp]
  \centering
  \caption{「华夏神族」名号攀附的三段流变}
  \label{tab:lineage}
  \footnotesize
  \begin{tabular}{@{}p{18mm}p{104mm}@{}}
    \toprule
    时期 & 使用状况与特征 \\
    \midrule
    秦汉 & 名号初起，范围甚窄，多见于尚存部分实传谱系与历数资格的小圈子 \\
    唐宋 & 谱牒之学大盛，攀附成风，名号正式与血统脱钩，成为通用郡望招牌 \\
    明清 & 余绪仍在；满清入主中原后以中华正统自居时，复将此旧称重新启用一次 \\
    \bottomrule
  \end{tabular}
  \tabsource{综合历代正史礼志、郡望谱牒及索科洛夫《记忆消除场的文献学后果》所引东亚谱系清退样本。}
\end{table}

然而到了清代以后，这么响亮、挂了两千多年的名号却在极短时间内销声匿迹，到如今只在少数冷僻笔记或伪史辨析里偶露一鳞半爪。把这一段陡然断掉的缘故摊开来看，实则分作前后相接的两层，先后次序不能颠倒。

第一层是实实在在的政治原因。清兵入关、鼎革既定之后，历法正朔另行颁布，明代以来的旧谱牒体制也随之废弛；「华夏神族」四个字不再能换来官位、婚嫁便利或者修史时的一席之地。一个招牌一旦彻底失去了用处，原本攥着它不放的那些家族和文人自己便首先闭了口，无意再去费神抄写或考辨它。

第二层才是前文反复碰到却始终没有深挖的那股推力：记忆消除场。当一群人因为现实无用而不再提起某个名号之后，它在日常判断里的分量就降到了零；一旦没有人再觉得它重要，剩下的工作便交给那个始终没有停止过运转的场，于是旧书不再重刻，家谱不再续修，老一辈临终也不再交代给子孙听。两层叠在一处，前后不过两三代人的工夫，一段两千年的称谓就干净地缩回了故纸堆里。第一层是门闩，第二层是门外的风；若是没有前面那一层政治上的「无用」，光靠后面那阵风，未必能把这件事抹得如此干净。

\section{末章的限度}

没有亲眼见过月面、只能在档案室里对着照片与残片拼图的笔者，把全书十二章的线头在这里收拢到一处，最诚实的做法依然是先交代哪些地方还站得不够稳。后半部年代愈晚、实物愈少，推论里掺进来的空白也就愈大。

这一章里直接用来定案的硬材料，严格说来只有第 117 全宗第 96 叶的存活率表与《广寒宫残简》简二十六，外加一份东渡名籍残档；至于公元前二一〇年那次跨海东渡的航程细节与列岛登岸经过，手上并没有同期日本本土留下的文字实录。其中最大的一处悬空，出在中原和日本之间那条跨海线头上。从公元前二一〇年到公元七一二年，中间隔着九百二十年几乎没有同期本地汉字记载的空白期；不仅如此，徐福船队的原始名籍早已一无所存。仅凭《史记》出航年份紧贴前二一一年东郡残片、编制又奇特地凑齐童男女与五谷百工，便断言八意永琳与蓬莱山辉夜带着另外二百一十五人就混在那几条船上，这一步同样是猜的。把「月读命」往回认成「月夜见」，把「八意思兼神」拆回「八意永琳」，音训与母题确实环环相扣，但音训扣得住并不等同于拿到了一份弥生时代的登岸花名册。

回头望去，这十二章书走过的路子越来越轻。头几章还有卫星轨道的二十九帧影像、贝加尔湖六点零三米深的孢粉层、四百一十二枚寸余厚的金属板压在桌面上，拿游标卡尺和同位素报告都量得出斤两；越往后翻，实物一片比一片少，年代一程比一程疏阔，到第十、十一章还有撤离签符、合金残片、弹道曲线与二里头烧结柱洞可资按图索骥，及至末章落入两千年的名号流变时，留下的便只剩几串读音和形貌还对得上齿的名字。那座城留在月面上的硬器物只够支撑到前二二一年断电那一刻；至于前二二一年以后在黄土和海岛上散开的人事，留在世上的也就只剩这一层稀薄的回声。
